\section{Considerações finais}\label{cap:fim}

\subsection{Síntese dos achados}

Esta revisão perguntou se o diagnóstico da literatura econômica sobre o impacto da inteligência artificial no mercado de trabalho foi alterado pelo avanço da IA generativa. A resposta é afirmativa, mas exige três qualificações que constituem, elas próprias, os principais resultados do trabalho.

A primeira diz respeito ao \emph{tamanho} da mudança. Entre os 816 estudos analisados, a fração que situa o risco nas ocupações de alta qualificação passa de 2,7\% entre os publicados até 2022 para 13,6\% entre os publicados de 2023 em diante, resultado que resiste a definições alternativas de período, à exclusão de subconjuntos do corpus e à correção para comparações múltiplas. O risco para a baixa qualificação, contudo, permanece o diagnóstico mais frequente nos dois períodos, o sinal predominante sobre o emprego continua ambíguo e o deslocamento segue sendo o mecanismo mais invocado. A hipótese de ruptura é confirmada em direção, mas não na forma forte de uma inversão do consenso.

A segunda diz respeito à \emph{origem} da mudança, e é o achado mais importante. O deslocamento do risco acompanha o objeto tecnológico, e não o calendário. Entre os estudos sobre IA generativa, o risco para a alta qualificação é a categoria modal (48,5\%); entre os estudos sobre robótica e automação, ele é residual e idêntico antes e depois do ChatGPT (1,8\% e 3,1\%). Usada como grupo de comparação, a literatura de robótica permite separar dois conjuntos de mudanças: as específicas do objeto, que são o deslocamento do risco e parte da alta da complementaridade, e as de época, que são o encurtamento do horizonte, o recuo da demanda agregada e o sinal menos negativo, presentes nas duas literaturas e atribuíveis ao amadurecimento empírico do campo. A leitura teórica é que a abordagem de tarefas sai confirmada: o risco segue o conteúdo das tarefas que cada tecnologia executa. O que envelheceu não foi a teoria, e sim a identificação entre tarefa rotineira e tarefa automatizável.

A terceira diz respeito à \emph{natureza} da evidência. O sinal que cada estudo atribui ao efeito sobre o emprego depende fortemente do desenho: estudos de exposição ocupacional tendem ao negativo, estudos de firma ao positivo, e modelos teóricos e revisões ao ambíguo. A aparente falta de convergência do campo é, em boa medida, um retrato da divisão de trabalho entre desenhos que capturam forças distintas do arcabouço de Acemoglu e Restrepo. Pela mesma razão, o novo diagnóstico sobre a IA generativa deve ser lido com cautela: os estudos que apontam a alta qualificação são exercícios de exposição, revisões e modelos, e nenhum é evidência de firma ou de nível agregado. A literatura de IA generativa encontra-se hoje onde a de robótica se encontrava há uma década, quando dispunha de índices de suscetibilidade e ainda não das estimativas causais que viriam a moderá-los.

\subsection{Limitações}

As conclusões estão sujeitas a limitações discutidas ao longo do trabalho e aqui reunidas. A busca tem sensibilidade limitada: recuperou 15 das 46 referências com data elegível entre as mais citadas pelo próprio corpus, não alcança \textit{working papers} e deixa de fora trabalhos de fronteira publicados em periódicos de ciência geral. O corpus representa, assim, a literatura econômica indexada e difusa, e não a fronteira da disciplina. Triagem, elegibilidade, extração e avaliação de qualidade foram realizadas por modelos de linguagem, e a verificação humana amostral prevista no protocolo é o instrumento que permite dimensionar o erro dessa cadeia. Apenas 14\% dos estudos foram lidos em texto completo, e a variável central, a polarização, é justamente a mais afetada: fica sem classificação em 40\% dos resumos e é simplificada quando classificada. As variáveis substantivas foram codificadas pelo mesmo modelo na mesma chamada, o que pode inflar associações entre elas; o teste da hipótese central, por apoiar-se no ano de publicação, é menos vulnerável a esse problema. O uso da robótica como grupo de comparação é uma análise \textit{post hoc}, e o grupo de IA generativa é pequeno. Por fim, os testes estatísticos são exploratórios: descrevem associações num conjunto de estudos, não efeitos numa população.

\subsection{Agenda de pesquisa futura}

Os resultados apontam quatro direções. A primeira é a produção de evidência causal sobre os efeitos da IA generativa no emprego, em nível de firma, de ocupação e de mercado de trabalho local, à semelhança do que a literatura de robótica construiu; sem ela, o diagnóstico de risco para a alta qualificação permanece uma hipótese sobre exposição. A segunda é a recuperação dos canais de equilíbrio geral: a literatura tornou-se mais empírica e mais microeconômica, mas a pergunta relevante para a política pública é agregada, e a demanda agregada é hoje o mecanismo menos estudado. A terceira é a distinção entre complementaridade e reinstalação: documenta-se cada vez mais que a IA eleva a produtividade nas tarefas existentes, e muito menos que cria tarefas novas, e é da segunda que depende a demanda por trabalho no longo prazo. A quarta é a dimensão intraocupacional: a evidência de compressão do retorno à experiência não cabe na dicotomia entre alta e baixa qualificação, e revisões futuras deveriam codificá-la separadamente. Do ponto de vista metodológico, a experiência desta revisão recomenda que campos deriváveis de metadados nunca sejam delegados ao extrator automático e que a sensibilidade da busca seja aferida contra um conjunto de referência antes da triagem.
